\documentclass[12pt]{article}
\usepackage[T1]{fontenc}
\usepackage[utf8]{inputenc}
\usepackage{lmodern}
\usepackage{textcomp}
\usepackage{enumitem}
\usepackage{geometry}
\geometry{margin=1in}
\begin{document}
\begin{center}
Answer Key - Philosophy - Undergraduate Level Assessment 1 \
\small (Answers are ordered exactly as in the question paper.)
\end{center}

\section*{Multiple Choice}
\begin{enumerate}[label=\arabic*.]
\item \textbf{Answer:} A
\\ \textit{Options:} A) both are infallible and beyond questioning.; B) ethics guides science in its pursuit of knowledge.; C) both are essential components of a worthwhile life.; D) both are completely separate and do not intersect.; E) both seek to provide us with true beliefs about the world.
\\ \textit{Note:} Hare argues that both ethics and science provide definitive truths that cannot be disputed, as they rely on rational foundations that aim for objective understanding and moral clarity.
\item \textbf{Answer:} D
\\ \textit{Options:} A) a detailed understanding of societal norms and expectations.; B) a comprehensive review of historical philosophical theories.; C) the influence of personal experiences on moral decisions.; D) analyses of concepts such as "action" and "intention."; E) the integration of religious beliefs into moral decisions.
\\ \textit{Note:} The correct answer highlights that Anscombe emphasizes the importance of understanding key concepts like "action" and "intention" in moral psychology, as they are fundamental to analyzing moral responsibility and ethical behavior.
\item \textbf{Answer:} A
\\ \textit{Options:} A) things that no one would want; B) what is considered deep insight; C) a life of poverty and hardship; D) the pressures of aristocratic society; E) what is considered utter despair
\\ \textit{Note:} Before his conversion, Tolstoy was surrounded by material possessions and societal values that he later came to view as empty and devoid of true meaning, leading him to reject them in search of a more authentic life.
\item \textbf{Answer:} A
\\ \textit{Options:} A) proportional; B) equality; C) retaliatory; D) punitive; E) compensatory
\\ \textit{Note:} Nathanson supports proportional retributivism because it asserts that the punishment should be commensurate with the severity of the crime, aligning with his belief in a fair and just response to wrongdoing.
\item \textbf{Answer:} D
\\ \textit{Options:} A) false cause; B) red herring; C) argumentum ad populum; D) jumping to a conclusion; E) ad novitatem
\\ \textit{Note:} The term "jumping to a conclusion" accurately describes a hasty conclusion because it refers to the process of making a judgment or decision based on insufficient evidence or reasoning, which is central to the concept of hasty generalization in philosophical discourse.
\end{enumerate}

\section*{True / False}
\begin{enumerate}[label=\arabic*.]
\setcounter{enumi}{5}
\item \textbf{Answer:} True
\\ \textit{Options:} A) True; B) False
\\ \textit{Note:} Lukianoff and Haidt argue that American colleges and universities promote critical reasoning by fostering an environment where students are encouraged to engage with diverse viewpoints and challenge their own beliefs through open discourse.
\item \textbf{Answer:} True
\\ \textit{Options:} A) True; B) False
\\ \textit{Note:} The statement is true because \_ad crumenam\_ is indeed a logical fallacy that occurs when one assumes that a person's wealth or financial success is indicative of their truthfulness or moral character, thereby conflating material wealth with virtue.
\item \textbf{Answer:} False
\\ \textit{Options:} A) True; B) False
\\ \textit{Note:} Hare actually characterizes individuals who prioritize ideals over people's interests as fanatics, rather than those who disregard ideals.
\item \textbf{Answer:} True
\\ \textit{Options:} A) True; B) False
\\ \textit{Note:} Stevenson argues that objective goodness cannot be fully comprehended without empirical experience and context, which means it cannot be understood a priori.
\item \textbf{Answer:} True
\\ \textit{Options:} A) True; B) False
\\ \textit{Note:} Moore's definition of the naturalistic fallacy highlights the error in equating moral properties, such as 'goodness,' with natural properties, asserting that moral qualities cannot be reduced to or identified with naturalistic descriptions.
\end{enumerate}

\section*{Long Form Response}
\begin{enumerate}[label=\arabic*.]
\setcounter{enumi}{10}
\item \textbf{Answer:} G.E. Moore's contributions to philosophy encompass several pivotal concepts, particularly his defense of common sense and the critique of idealism. One of his most significant ideas is the "open question argument," which challenges the equating of moral properties with natural properties, emphasizing that moral truths cannot simply be reduced to empirical observations. Additionally, Moore's notion of "good" as a non-natural property prompts ongoing debates about ethical realism and the nature of moral facts. The implications of these concepts are profound, as they challenge philosophers to reconsider the foundations of ethical theory and the relationship between language, meaning, and reality. Ultimately, Moore's work remains relevant as it encourages a rigorous examination of philosophical assumptions and the importance of clarity in ethical discussions.
\\ \textit{Note:} correct because it accurately identifies G.E. Moore's key philosophical contributions, such as the open question argument and the non-natural property of "good," and explains their significance in challenging moral reductionism and fostering debates in ethical realism.
\item \textbf{Answer:} The fallacy of hypostatization occurs when abstract concepts or terms are treated as if they have concrete existence, leading to misleading arguments. For instance, when one argues that "society demands change," the term "society" is often treated as a tangible entity rather than a complex network of individuals and relationships. This shift in meaning can obscure the fact that societal change involves specific actions by real people, thus weakening the argument's validity. Such misinterpretation can lead to oversimplified conclusions or unjustified assertions about collective behavior, ultimately undermining logical reasoning in philosophical discourse. Hypostatization illustrates the importance of clarity in language to maintain the integrity of an argument.
\\ \textit{Note:} correct because it illustrates how the fallacy of hypostatization misrepresents abstract terms as concrete entities, leading to flawed reasoning by oversimplifying complex social dynamics and obscuring the true nature of the argument.
\end{enumerate}

\vfill
\noindent\textit{End of Answer Key}
\end{document}